% !TEX root = main.tex

The development of natural user interfaces has fostered a growing interest in
hand gesture recognition as a means of interaction between humans and machines.
Gestures, being rapidly executed freehand drawings
\cite{Huang2019}, exhibit a strong inherent variability: two instances of the
same gesture produced by the same user may differ in speed, scale, spatial
position, and duration. The sequential nature of the signal and this variability
are the main challenges any recognizer must address.

In this work we consider the problem of classifying three-dimensional gesture
trajectories recorded as a sequence of spatial coordinates. Following the
project guidelines, we assume constant time steps ($t = 1, 2, 3, \dots$) and
ignore the precise sampling instants. Four classifiers are implemented and
compared:
\begin{enumerate}\itemsep 0em
    \item Dynamic Time Warping (DTW) as a first baseline, combined with
    a $k$-nearest-neighbors ($k$-NN) classifier;
    \item Edit-distance on quantized symbolic sequences, as a second
    baseline;
    \item The Three-Cents ($3\mathbb{c}$) recognizer, a 3D extension of the
    \$1 recognizer~\cite{Wobbrock2007};
    \item A Bidirectional Long Short-Term Memory network (BiLSTM) trained end-to-end.
\end{enumerate}

Each method is evaluated on two datasets from~\cite{Huang2019} under both a
\emph{user-independent} and a \emph{user-dependent} cross-validation protocol.
The research questions underlying this project are threefold: (Q1)~\emph{how do
classical template-matching baselines compare with more recent recognizers on
3D mid-air gestures?}; (Q2)~\emph{do user-independent settings systematically
degrade accuracy, as expected?}; and (Q3)~\emph{is one of the investigated
classifiers significantly better than the others in the user-independent
setting?}

The remainder of the report is organized as follows. Section~\ref{sec:data}
describes the datasets. Section~\ref{sec:theory} develops the theoretical
framework of the four methods. Section~\ref{sec:impl} outlines the
implementation and evaluation methodology. Section~\ref{sec:results} presents
and discusses the experimental results, and Section~\ref{sec:conclusion}
concludes.